% =============================================================================
%  1.1 Движок браузера
% =============================================================================
\section{Браузер Chromium. Структурная схема браузера и принцип работы компонентов движка}

\term{Chromium} - проект веб-браузера с открытым исходным кодом, на базе которого построены такие браузеры, как
Google Chrome, Microsoft Edge, Opera и другие. Структурная схема браузера изображена на рис.~\ref{fig:chromium-arch}.

\begin{figure}[H]
\centering
\begin{tikzpicture}[
  box/.style={font=\small, align=center, minimum height=12mm,
              minimum width=32mm, rounded corners=3pt,
              inner xsep=10pt, inner ysep=6pt, thick},
  parser/.style  ={box, fill=orange!15, draw=orange!50},
  tree/.style    ={box, fill=orange!8,  draw=orange!40},
  stage/.style   ={box, fill=orange!8,  draw=orange!40},
  jsblock/.style ={box, fill=teal!12,   draw=teal!40},
  gpu/.style     ={box, fill=green!10,  draw=green!40},
  screen/.style  ={box, fill=red!8,     draw=red!30},
  arr/.style={->, >={Stealth[length=6pt]}, line width=1.8pt, draw=gray!55},
  arrbind/.style={<->, >={Stealth[length=6pt]}, line width=1.8pt, draw=orange!70},
  scale=1.0, every node/.append style={transform shape},
]

% ==================== Контейнеры (рисуем первыми — фон) ====================

% Chromium
\fill[gray!6, rounded corners=8pt]  (-2, 2.4) rectangle (11.5, -12);
\draw[gray!40, very thick, rounded corners=8pt] (-2, 2.4) rectangle (11.5, -12);
\node[font=\normalsize\bfseries, anchor=south, transform shape=false]
     at (4.75, 2.4) {Chromium (браузер)};

% Blink
\fill[blue!5, rounded corners=6pt]  (-1.5, 1.4) rectangle (6.5, -9.8);
\draw[blue!30, thick, rounded corners=6pt] (-1.5, 1.4) rectangle (6.5, -9.8);
\node[font=\small\bfseries, color=blue!50, anchor=north, transform shape=false]
     at (2.5, 1.3) {Blink (движок рендеринга)};

% V8
\fill[teal!5, rounded corners=6pt]  (7, 1.4) rectangle (11, -3.8);
\draw[teal!30, thick, rounded corners=6pt] (7, 1.4) rectangle (11, -3.8);
\node[font=\small\bfseries, color=teal!60, anchor=north, transform shape=false]
     at (9, 1.3) {V8 (движок JS)};

% ==================== Узлы ====================

% Парсеры — CSS слева, HTML справа (чтобы DOM был рядом с V8!)
\node[parser] (cssp)  at (0.5, 0)    {CSS parser};
\node[parser] (htmlp) at (4.5, 0)    {HTML parser};
\node[jsblock](jsp)   at (9, 0)      {JS parser};

% Деревья — CSSOM слева, DOM справа (рядом с JIT!)
\node[tree]   (cssom) at (0.5, -2.2) {CSSOM};
\node[tree]   (dom)   at (4.5, -2.2) {DOM tree};
\node[jsblock](jit)   at (9, -2.2)   {JIT compiler};

% Pipeline (шаг 2.2 между центрами = одинаковые промежутки)
\node[stage]  (layout) at (2.5, -4.4)  {Layout tree};
\node[stage]  (paint)  at (2.5, -6.6)  {Paint};
\node[stage]  (comp)   at (2.5, -8.8)  {Compositing};

% GPU + Экран
\node[gpu]    (gpup) at (9, -8.8)    {GPU process};
\node[screen] (scrn) at (9, -11.0)   {Экран};

% ==================== Стрелки ====================

% Парсеры → деревья (вертикальные, чистые)
\draw[arr] (cssp)  -- (cssom);
\draw[arr] (htmlp) -- (dom);
\draw[arr] (jsp)   -- (jit);

% Деревья → Layout tree (сходятся в одну точку, одна стрелка вниз)
\coordinate (merge) at (2.5, -3.3);
\draw[line width=1.8pt, draw=gray!55] (cssom.south) |- (merge);
\draw[line width=1.8pt, draw=gray!55] (dom.south)   |- (merge);
\draw[arr] (merge) -- (layout.north);

% Pipeline вниз
\draw[arr] (layout) -- (paint);
\draw[arr] (paint)  -- (comp);

% Bindings: DOM ↔ JS parser — прямая диагональная линия
\draw[arrbind] (dom.north east) -- (jsp.south west)
  node[midway, above, sloped, font=\scriptsize, transform shape=false]{Bindings};

% Compositing → GPU process
\draw[arr] (comp.east) -- (gpup.west);

% GPU → Экран
\draw[arr] (gpup) -- (scrn);

\end{tikzpicture}
\caption{Структурная схема браузера Chromium}
\label{fig:chromium-arch}
\end{figure}